\documentclass[12pt]{article}
\usepackage{amsmath, amssymb}

\title{CSE400 --- Lecture L11\_12 \\ Transformation of Random Variables}
\author{Dhaval Patel, PhD}
\date{}

\begin{document}

\maketitle

\section{Definitions}

\textbf{Transformation of Random Variable}

Given:
\[
X = \text{Random Variable with known PDF } f_X(x)
\]
\[
Y = g(X)
\]

\section{Step-by-step Method (Monotonic Case)}

\textbf{Step 1: CDF of $Y$}
\[
F_Y(y) = Pr(Y \le y)
\]
\[
= Pr(g(X) \le y)
\]

If inverse exists:
\[
= Pr(X \le g^{-1}(y))
\]
\[
= F_X(g^{-1}(y))
\]

\textbf{Step 2: Differentiate to get PDF}
\[
f_Y(y) = \frac{d}{dy} F_Y(y)
\]
\[
= \frac{d}{dy}\left[F_X(g^{-1}(y))\right]
\]

Using chain rule:
\[
f_Y(y) = f_X(g^{-1}(y)) \cdot \frac{d}{dy} g^{-1}(y)
\]

\subsection{Case: Decreasing Function}

\[
F_Y(y) = Pr(Y \le y)
\]
\[
= Pr(X \ge g^{-1}(y))
\]
\[
= 1 - F_X(g^{-1}(y))
\]

Differentiate:
\[
f_Y(y) =
- f_X(g^{-1}(y)) \cdot \frac{d}{dy} g^{-1}(y)
\]

\section{Example}

Given:
\[
X \sim \text{Uniform}(-1,1)
\]

\[
f_X(x) =
\begin{cases}
\frac{1}{2}, & -1 < x < 1 \\
0, & \text{otherwise}
\end{cases}
\]

Find PDF of:
\[
Y = \sin\left(\frac{\pi X}{2}\right)
\]

\textbf{Solution}

Inverse:
\[
x = \frac{2}{\pi} \sin^{-1}(y)
\]

Derivative:
\[
\frac{dx}{dy} = \frac{2}{\pi} \cdot \frac{1}{\sqrt{1-y^2}}
\]

Using formula:
\[
f_Y(y) = f_X(x)\left|\frac{dx}{dy}\right|
\]

\[
= \frac{1}{2} \cdot \frac{2}{\pi} \cdot \frac{1}{\sqrt{1-y^2}}
\]

\[
= \frac{1}{\pi\sqrt{1-y^2}}, \quad -1<y<1
\]

\section{Function of Two Random Variables}

Let:
\[
Z = X + Y
\]

Find:
\[
f_Z(z)
\]

If $X$ and $Y$ independent, special distributions apply.

\section{CDF Method for Sum}

\[
F_Z(z) = Pr(Z \le z)
\]
\[
= Pr(X+Y \le z)
\]
\[
= \iint_{x+y \le z} f_{XY}(x,y)\,dx\,dy
\]

\section{Leibniz Rule}

Let:
\[
G(x) = \int_{a(x)}^{b(x)} h(x,y)\,dy
\]

Then:
\[
\frac{d}{dx}G(x)
=
\frac{db(x)}{dx}h(x,b(x))
-
\frac{da(x)}{dx}h(x,a(x))
+
\int_{a(x)}^{b(x)}
\frac{\partial}{\partial x}h(x,y)\,dy
\]

\section{Derivation of PDF of Sum}

\[
f_Z(z)
=
\int_{-\infty}^{\infty}
\left[
\frac{d}{dz}
\int_{-\infty}^{z-y}
f_{XY}(x,y)\,dx
\right] dy
\]

After applying Leibniz rule:
\[
f_Z(z)
=
\int_{-\infty}^{\infty}
f_{XY}(z-y,y)\,dy
\]

Alternative:
\[
f_Z(z)
=
\int_{-\infty}^{\infty}
f_{XY}(x,z-x)\,dx
\]

\section{Independent Case (Convolution)}

If independent:
\[
f_{XY}(x,y) = f_X(x)f_Y(y)
\]

Therefore:
\[
f_Z(z)
=
\int_{-\infty}^{\infty}
f_X(x)f_Y(z-x)\,dx
\]

Called \textbf{Convolution Integral}.

\section{Normal Example}

Given:
\[
X \sim N(0,1), \quad Y \sim N(0,1)
\]

\[
f_X(x) = \frac{1}{\sqrt{2\pi}} e^{-x^2/2}
\]
\[
f_Y(y) = \frac{1}{\sqrt{2\pi}} e^{-y^2/2}
\]

Using convolution:
\[
f_Z(z)
=
\int_{-\infty}^{\infty}
\frac{1}{\sqrt{2\pi}}e^{-x^2/2}
\cdot
\frac{1}{\sqrt{2\pi}}e^{-(z-x)^2/2}
\,dx
\]

Result:
\[
Z \sim N(0,2)
\]

\section{Exponential Example}

Given:
\[
f_X(x)=\lambda e^{-\lambda x}, \quad x>0
\]
\[
f_Y(y)=\lambda e^{-\lambda y}, \quad y>0
\]

Then:
\[
f_Z(z)
=
\int_0^z
\lambda e^{-\lambda x}
\cdot
\lambda e^{-\lambda(z-x)}
dx
\]

\[
= \lambda^2 e^{-\lambda z} \int_0^z dx
\]

\[
= \lambda^2 z e^{-\lambda z}, \quad z>0
\]

\section{Difference of Random Variables}

\[
Z = X-Y
\]

\[
F_Z(z)=Pr(X-Y \le z)
\]

\[
f_Z(z)=\frac{d}{dz}F_Z(z)
\]

Same procedure as sum.

\section{Ratio Transformation}

\[
Z = \frac{X}{Y}
\]

\[
F_Z(z)=Pr\left(\frac{X}{Y} \le z\right)
\]

Cases depend on:
\[
Y>0 \quad \text{or} \quad Y<0
\]

\section{Radial Transformation}

\[
R = X^2 + Y^2
\]

Find:
\[
f_R(r)
\]

\section{Important Results Summary}

\subsection{Single Variable Transformation}

\[
f_Y(y)=
f_X(g^{-1}(y))
\left|
\frac{d}{dy}g^{-1}(y)
\right|
\]

\subsection{Sum of Two RVs}

\[
f_Z(z)=\int f_{XY}(x,z-x)\,dx
\]

\subsection{Independent Case}

\[
f_Z(z)=\int f_X(x)f_Y(z-x)\,dx
\]

\subsection{Normal Sum Result}

\[
X,Y \sim N(0,1)
\Rightarrow
Z=X+Y \sim N(0,2)
\]

\subsection{Exponential Sum Result}

\[
f_Z(z)=\lambda^2 z e^{-\lambda z}, \quad z>0
\]

\end{document}
